\chapter{Introducción}

\section{Contexto y Planteamiento del Problema}

La electroencefalografía (EEG) constituye una de las herramientas más consolidadas en el ámbito de la neurociencia clínica y experimental. Su alta resolución temporal —del orden de milisegundos—, su carácter no invasivo y su costo operativo moderado en comparación con otras técnicas de neuroimagen, como la Imagen por Resonancia Magnética Funcional (fMRI) o la Tomografía por Emisión de Positrones (PET), la han posicionado como el método de referencia para el estudio de la dinámica cerebral en tiempo real. Sus aplicaciones abarcan desde el diagnóstico clínico de patologías neurológicas —epilepsia, trastornos del sueño, encefalopatías— hasta la investigación fundamental en neurociencia cognitiva y el desarrollo de interfaces cerebro-computador (BCI, \textit{Brain-Computer Interfaces}) para asistencia motora y comunicación aumentativa.

No obstante, la adquisición de señales EEG es un proceso notablemente frágil. En condiciones reales de laboratorio u hospital, es frecuente enfrentar la degradación o pérdida total de la señal en uno o más electrodos. Las causas son diversas: aumento progresivo de la impedancia de contacto por evaporación del gel conductor, fallas en los amplificadores diferenciales (saturación de voltaje), movimientos del paciente que desplazan físicamente los sensores, o simplemente el descarte deliberado de canales durante la etapa de preprocesamiento automático por presentar una relación señal-ruido inaceptable. En registros de alta densidad (32--256 canales), la probabilidad de que al menos un canal presente datos inválidos en una sesión típica de 30 minutos es cercana a la certeza.

La pérdida de canales no constituye un mero inconveniente cosmético: representa una amenaza concreta para la integridad del análisis neurofisiológico posterior. La ausencia de información espacial interrumpe la distribución topográfica del potencial eléctrico, distorsiona los algoritmos de localización de fuentes (\textit{source localization}), y altera biomarcadores clínicos de primer orden como los Potenciales Relacionados a Eventos (ERPs) y la Densidad Espectral de Potencia (PSD). En el contexto de los sistemas modernos de aprendizaje de máquina aplicados a BCI, la pérdida impredecible de canales modifica la dimensionalidad de los vectores de entrada, forzando el descarte de ensayos completos de datos cuya adquisición ha demandado tiempo y recursos considerables. En este escenario, la interpolación de canales faltantes deja de ser un paso opcional y se convierte en una necesidad metodológica insoslayable dentro del \textit{pipeline} de preprocesamiento.

\section{Antecedentes Específicos y Estado del Arte}

El problema de la interpolación de canales EEG ha recibido atención sostenida desde la década de 1980, generando una familia diversa de soluciones que pueden agruparse en cuatro grandes paradigmas.

\subsection{Paradigma Geométrico-Espacial}

La estrategia más extendida, implementada como opción por defecto en las principales plataformas de análisis —EEGLAB~\cite{delorme2011eeglab}, MNE-Python~\cite{gramfort2013meg}, Brainstorm~\cite{taadel2011brainstorm}— es la interpolación por \textit{Spherical Splines}, introducida por Perrin et al.~\cite{perrin1989}. Este método modela la cabeza como una esfera conductora y representa el potencial eléctrico mediante una expansión en polinomios de Legendre. Su solidez numérica y su costo computacional reducido lo han mantenido como estándar \textit{de facto} durante más de tres décadas. Variantes posteriores incluyen funciones de base radial (RBF) con kernel de placa delgada (\textit{Thin-Plate Spline}) y el ponderamiento por distancia inversa (IDW). La limitación transversal de esta familia es que cada muestra temporal se interpola de forma aislada, sin memoria del contexto dinámico de la señal.

\subsection{Paradigma Estadístico}

Métodos como el Análisis de Componentes Independientes (ICA) proyectan las señales observadas en un espacio de fuentes latentes estadísticamente independientes, reconstruyendo los canales faltantes desde la proyección inversa~\cite{jung2000removing}. Aunque efectivos para remover artefactos oculares y musculares, su desempeño en interpolación se degrada marcadamente bajo pérdida severa, especialmente cuando los canales afectados son espacialmente contiguos.

\subsection{Paradigma de Aprendizaje Profundo (Deep Learning)}

En años recientes, los métodos basados en redes neuronales han comenzado a explorarse para la reconstrucción de señales EEG. Autoencoders convolucionales~\cite{banville2021uncovering} aprenden representaciones comprimidas de segmentos EEG multivariados y reconstruyen canales faltantes desde el espacio latente. Redes generativas antagónicas (GANs) han sido propuestas para imputación de series temporales neurofisiológicas, y arquitecturas de grafos neuronales (GNNs)~\cite{defferrard2016} explotan directamente la estructura de conectividad entre electrodos. Si bien estos enfoques pueden alcanzar desempeño competitivo en dominios restringidos, presentan limitaciones significativas que motivan su exclusión del presente estudio: (a) requieren grandes volúmenes de datos etiquetados para entrenamiento supervisado, (b) son vulnerables al desplazamiento de dominio (\textit{domain shift}) entre sujetos, paradigmas experimentales y configuraciones de electrodos, y (c) carecen de garantías formales de convergencia o de interpretabilidad sobre el mecanismo de reconstrucción.

En contraste, los métodos algorítmicos basados en optimización —como los que aborda esta tesis— son \textit{training-free}: operan directamente sobre la señal de prueba utilizando únicamente la topología de electrodos y \textit{priors} físicamente motivados de suavidad espacial y temporal. Esta propiedad los hace inmediatamente aplicables a cualquier montaje EEG sin requerir fase de entrenamiento, y sus soluciones están respaldadas por garantías matemáticas de optimalidad.

\subsection{Paradigma de Procesamiento de Señales en Grafos (GSP)}

En la última década, el marco teórico del GSP~\cite{shuman2013,ortega2018} ha ofrecido una alternativa conceptualmente superior a los enfoques puramente espaciales. En lugar de asumir un modelo esférico continuo, GSP representa la red de electrodos como un grafo discreto $\mathcal{G} = (\mathcal{V}, \mathcal{E}, \mathbf{W})$, donde cada electrodo es un vértice y los pesos $W_{ij}$ codifican relaciones de proximidad física o acoplamiento funcional entre canales. Bajo este paradigma, interpolar un canal faltante equivale a resolver un problema de optimización con restricciones de suavidad sobre la topología del grafo~\cite{narang2013,puy2018}.

Extensiones recientes incorporan además regularización temporal —esto es, penalizan cambios abruptos entre muestras consecutivas— a través de formulaciones como la Regularización Temporal de Suavidad de Sobolev (TRSS)~\cite{giraldo2022}, que extiende el funcional de costo de Tikhonov clásico con un operador de diferencias finitas temporales que refleja matemáticamente la inercia de los generadores neuronales. Otros trabajos han explorado variantes como el Laplaciano temporal (producto de Kronecker entre Laplacianos espaciales y temporales), la difusión por kernel de calor, y la construcción de grafos mediante grafos de visibilidad con kernel no-negativo (NNK)~\cite{bozkurt2022visibility}.

\section{La Brecha Metodológica y la Motivación de Este Trabajo}

Pese a la madurez de estos desarrollos, la literatura especializada exhibe dos déficits metodológicos que comprometen la comparabilidad de los resultados publicados.

En primer lugar, los estudios que proponen nuevos algoritmos de interpolación suelen evaluarlos contra \textit{baselines} ejecutados con sus parámetros por defecto —los que vienen preconfigurados en las bibliotecas de software—, mientras que el método propuesto es cuidadosamente sintonizado. Esta asimetría infla artificialmente la ventaja reportada y oscurece el verdadero techo de rendimiento de los métodos establecidos~\cite{calazans_machine_2024}.

En segundo lugar, la evaluación se restringe típicamente a una o dos métricas de error de amplitud global —MAE o RMSE— sobre uno o dos datasets, frecuentemente con simulaciones de pérdida de canales que no reflejan los modos de falla observados en la práctica clínica: la pérdida contigua de múltiples canales por desprendimiento de un parche del gorro EEG o por falla de un módulo de amplificación local. Muy pocos trabajos examinan el impacto sobre la morfología temporal (ERPs) o la preservación del contenido espectral (PSD, coherencia), que son, en última instancia, los predictores de utilidad clínica y científica de la señal interpolada.

La motivación central de esta tesis emerge directamente de estas brechas. El objetivo no es proponer un nuevo algoritmo de interpolación, sino someter a la frontera del conocimiento actual —métodos geométricos clásicos, GSP espacial y GSP con regularización temporal— a la comparación empírica más rigurosa, exhaustiva y justa de la que se tenga registro.

\section{Desarrollo Cronológico del Proyecto}

Dado que esta investigación constituye un esfuerzo de ingeniería experimental de gran escala, es pertinente describir la secuencia metodológica que condujo desde las decisiones iniciales de diseño hasta los resultados finales. Esta cronología no solo documenta el proceso, sino que transparenta las iteraciones de refinamiento que son inherentes a todo \textit{benchmark} computacional ambicioso.

\subsection{Fase 1: Selección y Estudio Bibliográfico de Métodos}

El punto de partida fue una revisión sistemática de la literatura en GSP aplicado a señales biomédicas. Se identificaron tres familias de métodos —geométricos, GSP instantáneos y GSP con regularización temporal— y se seleccionaron 18 algoritmos representativos de cada familia. Esta selección no fue exhaustiva sino \textit{representativa}: se priorizaron métodos con formulaciones matemáticas distintas (splines, Tikhonov, BGSRP, TRSS, Sobolev, Laplaciano temporal, difusión térmica), de modo que el \textit{benchmark} cubriera un espectro amplio de estrategias de regularización.

\subsection{Fase 2: Implementación del Pipeline Modular en Python}

Se desarrolló un \textit{pipeline} computacional estructurado en cinco módulos independientes: carga de datos (con MNE-Python como \textit{backend}), construcción de grafos (8 constructores), interpolación (18 métodos), evaluación (6 métricas) y orquestación. Cada módulo expone interfaces homogéneas, lo que permite intercambiar componentes sin modificar el resto del sistema. La implementación de TRSS siguió fielmente la formulación de Giraldo et al.~\cite{giraldo2022}, con mejoras ingenieriles como paso de gradiente adaptativo basado en la cota de Lipschitz del operador compuesto, \textit{clipping} por escala de datos, e inicialización por media de canal en lugar de ceros.

\subsection{Fase 3: Experimentación Exploratoria (Screening)}

Con los 18 métodos implementados y los 8 constructores de grafo, se ejecutó una fase exploratoria inicial sobre datos sintéticos. El objetivo no era obtener resultados definitivos, sino identificar qué métodos y constructores mostraban estabilidad numérica y desempeño competitivo, y cuáles presentaban fallas sistemáticas (singularidad del Laplaciano, no-convergencia, sensibilidad extrema a hiperparámetros). Esta fase permitió descartar aproximadamente dos tercios del espacio de métodos y reducir el foco a las configuraciones con mayor potencial.

\subsection{Fase 4: Selección de Métodos para Optimización Fina}

Del \textit{screening} emergió un subconjunto de 5 métodos que combinaban desempeño competitivo con estabilidad numérica: TRSS (mejor desempeño global en datos sintéticos), Tikhonov en grafo (mejor método GSP instantáneo), Spherical Spline y RBF Thin-Plate (mejores baselines geométricos), e ICA (mejor baseline estadístico). Para la construcción de grafos, el $k$-NN Gaussiano se seleccionó como constructor primario por su balance óptimo entre conectividad garantizada, eficiencia computacional y desempeño. Los métodos de grafo avanzados (Kalofolias, NNK, AEW) y las variantes temporales alternativas (Laplaciano temporal, difusión térmica) se documentaron pero no se incluyeron en la fase de optimización exhaustiva, dado que su costo computacional era desproporcionado frente a la ganancia marginal observada.

\subsection{Fase 5: Optimización Bayesiana con Optuna}

La fase central del proyecto consistió en someter los 5 métodos seleccionados a optimización bayesiana mediante Optuna~\cite{akiba2019optuna} sobre los tres datasets reales. Se emplearon 30 \textit{trials} por combinación (dataset $\times$ escenario $\times$ método) para los métodos GSP y 10 \textit{trials} para los baselines. Para prevenir fuga de datos en señales con alta autocorrelación temporal, se implementó un protocolo de validación con separación cronológica estricta: 70\% de cada segmento para optimización, un \textit{buffer} de 2 segundos descartado, y el remanente como conjunto de prueba independiente. Todas las métricas reportadas en esta tesis corresponden exclusivamente al conjunto de prueba.

\subsection{Fase 6: Comparación Confirmatoria contra MNE-Python}

Una vez caracterizado el desempeño de TRSS dentro del \textit{pipeline} propio, se reconoció la necesidad de una validación externa contra la implementación de referencia que cualquier investigador en EEG utilizaría en la práctica: \texttt{mne.io.Raw.interpolate\_bads(method='spline')} de MNE-Python. Esta comparación se diseñó como un experimento confirmatorio separado, con su propio protocolo de máscaras estáticas, optimización independiente de TRSS, y evaluación multi-métrica. La Sección~\ref{sec:mne_confirmatorio} del Capítulo~2 y la Sección~4.7 del Capítulo~4 documentan este \textit{benchmark} en detalle.

\subsection{Fase 7: Ablación del Componente Temporal}

Como validación adicional de la hipótesis, se ejecutó un estudio de ablación que compara TRSS-Full ($\alpha > 0, \beta > 0$) contra TRSS-NoTemporal ($\alpha > 0, \beta = 0$) y contra el baseline de Spline Espacial. Este experimento aísla la contribución marginal del término de regularización temporal y cuantifica, con pruebas de significancia estadística (Mann-Whitney U, delta de Cliff con corrección FDR), en qué métricas y bajo qué condiciones el término $\beta$ aporta valor.

\section{Objetivos}

\subsection{Objetivo General}

Comparar el rendimiento espacial, temporal y espectral de los métodos de reconstrucción de canales EEG faltantes basados en Procesamiento de Señales en Grafos (GSP) y regularización temporal frente a \textit{baselines} geométricos clásicos, mediante un marco experimental riguroso, reproducible y con optimización imparcial de hiperparámetros.

\subsection{Objetivos Específicos}

\begin{enumerate}
    \item \textbf{Unificación metodológica.} Implementar y consolidar en un \textit{pipeline} modular en Python una colección amplia de métodos de interpolación —geométricos, estadísticos y basados en grafos con regularización temporal—, garantizando interfaces homogéneas y evaluación estandarizada.

    \item \textbf{Diseño de escenarios clínicamente representativos.} Definir protocolos de simulación de canales faltantes que abarquen desde fallos esporádicos independientes (\textit{random}) hasta pérdidas masivas y espacialmente contiguas (\textit{nearby}), con niveles progresivos de severidad desde el 10\% hasta el 40\% de los canales totales.

    \item \textbf{Optimización imparcial (\textit{Fair Benchmarking}).} Integrar el \textit{framework} Optuna con muestreadores TPE y NSGA-II para explorar los espacios paramétricos de \textit{todos} los métodos —incluyendo los \textit{baselines} clásicos—, eliminando el sesgo de comparación contra configuraciones por defecto que permea la literatura.

    \item \textbf{Evaluación multi-métrica y multi-dominio.} Cuantificar el desempeño de reconstrucción sobre tres bases de datos de distinta densidad espacial y paradigma experimental (PhysioNet EEGMMIDB, BCI Competition IV 2a, MNE Sample), utilizando seis métricas complementarias: MAE, RMSE, SNR, DTW, LSD y Coherencia.

    \item \textbf{Análisis de relevancia fisiológica.} Demostrar cualitativa y cuantitativamente el impacto de cada familia de métodos sobre la preservación de Potenciales Relacionados a Eventos (N100, P300) y de la densidad espectral en las bandas de interés clínico (Delta, Theta, Alfa, Beta, Gamma).

    \item \textbf{Validación confirmatoria contra el estándar clínico.} Comparar el desempeño del mejor método GSP-temporal (TRSS) contra la implementación de referencia \texttt{MNE interpolate\_bads} bajo condiciones controladas de máscara estática y optimización independiente, estableciendo recomendaciones condicionales basadas en evidencia.

    \item \textbf{Trazabilidad y reproducibilidad.} Generar un repositorio de artefactos computacionales —configuraciones, resultados numéricos, bitácoras de ejecución y advertencias— que permita la auditoría completa de cada resultado reportado y su reproducción independiente.
\end{enumerate}

\section{Hipótesis de Trabajo}

La hipótesis que guía esta investigación es la siguiente:

\begin{quote}
La incorporación de restricciones de suavidad espacial modeladas mediante topologías de grafo, en conjunto con \textit{priors} de continuidad temporal penalizada —tal como los que provee la regularización de Sobolev TRSS—, permite reconstruir canales EEG faltantes con una precisión significativamente superior a la de los métodos geométricos clásicos de extrapolación de superficies. Esta superioridad se manifiesta no solo en la reducción del error de amplitud, sino fundamentalmente en la preservación de la morfología de los potenciales evocados (ERPs) y de la distribución de potencia espectral (PSD), y se acentúa bajo condiciones de pérdida espacial severa y contigua.
\end{quote}

Al cierre de la fase experimental de esta investigación —que abarcó cientos de iteraciones sistemáticas y decenas de miles de evaluaciones de interpolación—, la evidencia acumulada es ampliamente favorable a esta hipótesis. Los resultados que se presentan en los Capítulos~4 y~5 muestran que el término de regularización temporal constituye un punto de inflexión en el desempeño de los métodos de reconstrucción, particularmente en los regímenes donde la información espacial se toma crítica. No obstante, la comparación confirmatoria contra MNE-Python introduce matices importantes: el \textit{spherical spline} industrializado de MNE es un \textit{baseline} excepcionalmente fuerte que, bajo condiciones favorables (montajes densos, pocos canales perdidos, pérdida aleatoria), ofrece desempeño competitivo con costo computacional drásticamente menor.

\section{Alcance y Contribuciones}

El alcance de este trabajo abarca la ejecución computacional a gran escala de un \textit{pipeline} de evaluación comparativa que cruza sistemáticamente tres bases de datos, dos modos de pérdida, cuatro niveles de severidad, múltiples constructores de grafo y varias familias de métodos de interpolación. En total, el espacio experimental comprende del orden de decenas de miles de configuraciones evaluadas bajo estricto control de semillas y con separación cronológica estricta entre datos de optimización y de prueba.

Las contribuciones principales de esta tesis se articulan en cuatro ejes:

\begin{enumerate}
    \item \textbf{Evidencia estadística del cruce de desempeño.} Se demuestra que, bajo condiciones de sintonización equitativa, la familia de métodos con regularización temporal —encabezada por TRSS— supera de manera consistente a los \textit{baselines} geométricos en las dimensiones de error de amplitud (MAE, RMSE), fidelidad morfológica (DTW) y preservación espectral (LSD, coherencia), con márgenes que se amplifican al aumentar la severidad de la pérdida.

    \item \textbf{Caracterización del colapso de los métodos clásicos.} Se documenta empíricamente que, bajo regímenes de pérdida contigua severa (40\% de los canales en una misma región del cuero cabelludo), los métodos puramente espaciales —incluyendo el \textit{Spherical Spline}— experimentan un colapso de fase en el que la señal reconstruida invierte su polaridad o diverge en amplitud, mientras que los métodos con anclaje temporal mantienen trayectorias fisiológicamente plausibles.

    \item \textbf{Validación confirmatoria contra el estándar de la industria.} Se presenta una comparación directa y rigurosa entre TRSS y \texttt{MNE interpolate\_bads}, documentando los escenarios donde cada enfoque es preferible. La conclusión no es de dominancia universal sino condicional: TRSS se justifica cuando la prioridad es la preservación de dinámica temporal fina bajo pérdida severa; MNE es la opción práctica correcta en escenarios de rutina con montajes densos.

    \item \textbf{Infraestructura de código abierto para la comunidad.} Se entrega un \textit{toolkit} de experimentación unificado en Python, modular, completamente rastreado mediante control de versiones y diseñado para que cualquier investigador pueda reproducir los resultados aquí presentados o extender el \textit{framework} con nuevos métodos, datasets o métricas.
\end{enumerate}

\section{Estructura del Documento}

El resto de esta tesis se organiza de la siguiente manera:

\begin{itemize}
    \item \textbf{Capítulo~2 -- Marco Teórico.} Presenta los fundamentos neurofisiológicos del EEG, la teoría matemática del Procesamiento de Señales en Grafos, las formulaciones de los métodos de interpolación clásicos y basados en GSP, un análisis detallado de la implementación de \texttt{MNE interpolate\_bads} como \textit{baseline} de referencia, y la definición formal de las métricas de evaluación.

    \item \textbf{Capítulo~3 -- Metodología.} Describe el diseño experimental completo: los tres datasets, el \textit{pipeline} computacional, los métodos de construcción de grafos, los protocolos de simulación de canales faltantes, la arquitectura de optimización bayesiana, el protocolo de validación temporal, y la infraestructura de evaluación estadística.

    \item \textbf{Capítulo~4 -- Experimentos y Resultados.} Expone los hallazgos empíricos organizados en niveles crecientes de especificidad: comparación global de familias, tablas estratificadas por severidad y modo de pérdida, análisis de preservación de ERPs y PSD, curvas de degradación, estudio de ablación del término temporal, y \textit{benchmark} confirmatorio contra MNE-Python.

    \item \textbf{Capítulo~5 -- Discusión.} Interpreta los resultados a la luz de la hipótesis de trabajo, analiza las causas metodológicas de las diferencias observadas, discute las implicaciones para la práctica clínica, y transparenta las limitaciones del estudio.

    \item \textbf{Capítulo~6 -- Conclusiones y Trabajo Futuro.} Sintetiza los hallazgos principales, evalúa el grado de cumplimiento de los objetivos y la hipótesis, y propone direcciones de investigación futura.
\end{itemize}